\section{Related work}
\label{sec:related}

\subsection{Knowledge management and document acquisition}

Recent work on knowledge-based systems combines language models with curated
knowledge \citep{yang-kbs-llm-2025}. DBpedia publishes structured
information from Wikipedia as linked data \citep{auer-dbpedia-kb-2007}.
Enterprise content management governs the documents through which
organisations retain and share knowledge
\citep{paivarinta-ecm-2005,tyrvaainen-ecm-2006}. Knowledge creation theory
also examines how organisations articulate tacit knowledge
\citep{nonaka-seci-1994}.
Metadata, tagging and human review affect which materials users can find and
trust \citep{park-metadata-quality-2009,golder-tagging-2006,mosqueira-rey-hitl-2022}.
Text categorisation supports placement \citep{sebastiani2002machine}.
Layout-aware parsing and MinerU recover structured content for indexing
\citep{zhong-publaynet-layout-2019,huang-layoutlmv3-docai-2022,mineru}, while
near-duplicate detection and provenance help maintain the collection
\citep{manku-near-dup-2007,missier-prov-2013,cheney-provenance-2009}.
HMRS connects these operations by making the results of acquisition and
placement visible to the workflow that selects evidence.

Gallofr\'e Oca\~na and Opdahl \citep{ocaa2023software} describe how to
integrate AI and knowledge bases in a newsroom platform. Related systems
ground dialogue in a graph, build open-domain question answering for
low-resource languages, or support retrieval by bank supervisors
\citep{yang2024building,budur2024building,aarab2024system}. LLMs can also
populate domain ontologies \citep{ciatto2025large}. HMRS applies this systems
perspective to changing document collections. Its focus is the managed
state a retrieval workflow reads and an operator can inspect after a
lifecycle operation.

\subsection{Retrieval, organisation and agentic search}

\begin{table*}[!t]
\centering
\small
\setlength{\tabcolsep}{3.5pt}
\renewcommand{\arraystretch}{1.02}
\caption{Representative systems compared along the dimensions examined here.
A dash indicates a dimension outside the cited work's main contribution.}
\label{tab:related-comparison}
\begin{tabular}{@{}>{\raggedright\arraybackslash}p{0.13\textwidth}>{\raggedright\arraybackslash}p{0.245\textwidth}>{\raggedright\arraybackslash}p{0.21\textwidth}>{\raggedright\arraybackslash}p{0.215\textwidth}>{\raggedright\arraybackslash}p{0.13\textwidth}@{}}
\toprule
System & Hierarchy and scope & Construction and lifecycle & Evidence handling & Access surface\\
\midrule
RAPTOR \citep{raptor} & cluster tree over document text & built offline at index time & --- & tree retrieval; prototype\\
GraphRAG \citep{graphrag} & entity--relation graph, community summaries & pre-generated offline & --- & graph retrieval; prototype\\
LightRAG \citep{lightrag} & low- and high-level graph retrieval & incremental graph updates & --- & graph retrieval; prototype\\
DeepRead \citep{deepread} & per-document structured Markdown, section coordinates & structure built at ingestion & agent locates then reads sections & agentic search system\\
EPRAG \citep{kinalioglu-kbs-eprag-2026} & multi-source enterprise corpus & --- & evidence-state diagnosis; abstention or clarification & enterprise retrieval\\
HMRS (this work) & managed collection--document--segment state & content-based placement; lifecycle operations; three-source reconciliation & optional gate, separate from ordering & self-hosted; MCP layer\\
\bottomrule
\end{tabular}
\end{table*}

RAG uses retrieved passages to inform generation
\citep{rag-lewis,gao-rag-survey-2023}. Dense
retrieval ranks passages in a learned representation space \citep{dpr}; BM25
uses lexical signals \citep{bm25}. Enterprise architectures combine semantic
retrieval with vector storage and LLMs \citep{ghali-kbs-retrieval-2025}.
Hybrid search combines ranking signals, for example with reciprocal-rank
fusion \citep{rrf-cormack}, and BGE-M3 provides multilingual and
multi-granularity embeddings \citep{bge-m3}. Structured retrieval organises
the units a ranker searches. GraphRAG builds entity graphs and community
summaries \citep{graphrag}. LightRAG searches an incrementally updated
graph at two retrieval levels \citep{lightrag}, and RAPTOR recursively
groups text into a cluster tree \citep{raptor}. RMR-SST represents multimodal engineering
documents as semantic trees \citep{li-kbs-rmrsst-2026}. Dense X Retrieval
indexes propositions \citep{chen2024dense}, while TEIE-RAG combines
hierarchical semantic chunking with source highlighting for engineering
information extraction \citep{shen-kbs-teierag-2026}.
Table~\ref{tab:related-comparison} locates HMRS among these designs. Its
hierarchy represents collection ownership and document structure, so a
query reads a maintained record of which collections and documents exist.
Index maintenance complements this hierarchy by keeping passages reachable
after document changes \citep{xu2023spfresh}.

Distributed IR treats collection selection and result merging as its two
subproblems \citep{callan1995searching,si2003redde,shokouhi-federated-2011}.
HMRS applies the same select-then-search principle to collection and
document descriptions that operators can maintain: the catalogue-first
path selects collections from their descriptions, and document selection
plays the role of result merging. Tool-using agents can choose and iteratively refine
what they retrieve \citep{react-yao,toolformer}; MCP provides standard
interfaces to these tools \citep{mcp}. Modular RAG studies how these
workflows combine retrieval components \citep{modular-rag}. DeepRead
illustrates structure-aware agent access by indexing paragraphs with section
coordinates and reading selected sections \citep{deepread}. HMRS gives
an agent a catalogue and shelf to determine which documents to open, then
exposes document and passage access through the same tool layer.

In engineering design, structured facts extracted from patents support RAG
\citep{siddharth2024retrieval}. Other approaches change which evidence a
query can reach through soft-prompt tuning, semantic query expansion and
stateful graph exploration
\citep{peng2024soft,massai2023evaluatio,tao2025fine}. HMRS exposes the document
state on which collection and passage selection depend.

\subsection{Verification, correction and abstention}

MemGPT and A-MEM maintain agent interaction records and notes
\citep{memgpt,amem}. HMRS maintains a corpus of source
documents and the descriptions, ownership and indices used to find them.
Knowledge organisation standards such as SKOS \citep{skos} and Dublin Core
\citep{dcmi-terms}
give conventions for vocabularies and descriptive metadata. At ingestion,
the HMRS workflow derives descriptions and tags from document content.

Self-RAG and CRAG assess evidence or correct retrieval during generation
\citep{selfrag,crag}; FLARE triggers retrieval as more evidence is needed
\citep{flare}. Work on selective prediction under domain shift and
uncertainty reporting examines when to withhold an answer
\citep{kamath,feng-abstain,yin-knowwhat}. Citation studies examine whether
references support the statements in the answer
\citep{liu-verifiability,gao-citations}.
EPRAG diagnoses evidence states to choose among response, clarification and
abstention in enterprise retrieval \citep{kinalioglu-kbs-eprag-2026}.
HMRS makes verification a server-side operation. The gate fetches
candidate text, scores passages and returns surviving text with provenance,
so an agent can inspect the source material behind the result.

Judge studies report preference bias, variable agreement and ranking
instability \citep{zheng-llmjudge,sun-urara,panickssery-selfpref}.
Logic-grounded metamorphic testing checks whether model decisions remain
consistent under equivalent input transformations \citep{zhou2026lgmt}.
HMRS tests evidence admission, candidate order and answer-stage reporting
as distinct decisions.

\subsection{Research gap}
\label{sec:gap}

RMR-SST and TEIE-RAG structure engineering document content for retrieval
and evidence access \citep{li-kbs-rmrsst-2026,shen-kbs-teierag-2026}. EPRAG
selects a response action once the candidate evidence is available
\citep{kinalioglu-kbs-eprag-2026}. Archival curation and enterprise
content management do govern placement and description over time
\citep{paivarinta-ecm-2005,park-metadata-quality-2009,missier-prov-2013};
the unstudied step is the coupling of that managed state to evidence
access, that is, whether a failed search can be traced to the placement,
description or index decision that supplied its candidates. HMRS studies the collection state that
supplies evidence to such workflows (Table~\ref{tab:related-comparison}
and Supplementary Table~S11). It records where documents belong, exposes
the descriptions used to select them, and lets operators check whether
the index reflects lifecycle changes. The evaluation follows this path from
post-ingest state and catalogue visibility through candidate selection and
verifier ordering.
